\section{Related Work and Scope}\label{sec:related}
\subsection{Sensitivity analysis and fault-injection prioritization}
FKeras analyzes the sensitivity of heavily quantized edge networks using gradients and leading weight-loss Hessian eigenpairs~\cite{weng2024fkeras}. It evaluates single- and multi-bit faults and refines parameter ordering within significance tiers. The present comparison preserves the published projection and tier ordering while adapting the continuous model and numerical solver to the frozen INT8 benchmarks. We do not equate FKeras with a diagonal-Hessian approximation or claim that it ignores cross-parameter curvature.

PrioriFI extends this line of work by using observed damage to guide subsequent fault injections~\cite{weng2026priorifi}. Hessian information initializes significance queues, and feedback determines which queue to probe next. Its endpoint is efficient discovery of sensitive fault locations. Converting that procedure into a mask of exactly $B$ protected bits introduces a further decision: how to combine observed and unobserved coordinates, how much probing to spend, and how to freeze the result before independent evaluation. NetSuRF has engineering support for studying such a conversion, but the present manuscript contains no completed PrioriFI-derived effectiveness comparison.

EISFINN uses importance sampling guided by neuron- and bit-level attribution to improve the efficiency of fault-injection campaigns~\cite{colucci2024eisfinn}. This is relevant precedent for directed exploration of faults, including an activation-level fault domain. Finding harmful events efficiently and estimating risk under a prescribed random law are nevertheless different statistical objectives. An importance-biased sample requires the corresponding probability correction before it estimates the latter. Our final evaluations retain the declared independent Bernoulli target and paired field identities.

These studies already establish that bit significance alone is an incomplete sensitivity description and that adaptive injection can be useful. Our empirical question is more specific: for frozen equal-budget masks, how does measured task risk vary with BER and coverage, and does exact conditioning improve a local protection decision beyond exact isolated measurements? The new measurements and numerical contracts, rather than the existence of nonlinearity or adaptive exploration, form the basis of the comparison.

\subsection{Selective hardening and physical mitigation}
Selective hardening has been investigated for FPGA neural networks~\cite{libano2019selective}; selective TMR can triplicate sensitive channels to meet a minimum-accuracy requirement under single channel-stuck faults~\cite{garciabertoa2023selective}. Other techniques exploit bit-level redundancy for recovery~\cite{catalan2025redundancy}. FORTUNE combines post-training affine quantization with two additional copies of each weight's MSB and majority voting during inference~\cite{nazari2024fortune}. Its precision search trades clean accuracy against reliability and resource requirements. This shares a significance preference with MSB-first allocation, but changes the representation and includes redundant storage and voting. Our fixed-checkpoint masks do not reproduce those correction mechanisms or their costs.

In particular, ideally protecting all INT8 sign bits covers 12.5\% of our eligible population, but implements neither ECC nor TMR. Ranger restricts intermediate value ranges to dampen transient datapath corruption, with protected memory assumed in its main fault model~\cite{chen2021ranger}. FT-CNN uses convolution checksums with distinct detection and correction capabilities; it also considers memory corruption and recovery through weight reloading or recomputation~\cite{zhao2021ftcnn}. These are complementary interventions, not interchangeable ranking methods. Their quantitative comparison would require matched fault domains, executable recovery behavior and verified resource costs. We therefore discuss them as related protection mechanisms without assigning their names to ideal-suppression curves.

The public \emph{Learned or Ranked?} companion release is another close precedent~\cite{alobaidi2026selectiveecc}. Its archived implementation includes matched-cost protection methods, held-out evaluation, paired layer interaction diagnostics and context-dependent layer marginal rankings. The release is cited as a dataset and source archive; peer-reviewed publication of its accompanying manuscript has not been established here. Its layer/bit-plane decisions and ideal or word-level ECC models differ from our individual-bit, fixed-cardinality interventions and reference integer execution. We do not claim the first observation of context-dependent protection, nor superiority to its allocation algorithms.

\subsection{Relationship to the earlier NetSuRF design flow}
Our FCCM work already connected susceptibility analysis, user budgets, selective TMR and an HLS-to-hardware implementation flow~\cite{guo2025inpixel}. The overall engineering motivation and the idea of using bit priorities to allocate resources are therefore inherited. This paper contributes a separate software methodology with exact encoded-state injection, explicit score-to-mask contracts, paired uncertainty, independent task-specific benchmarks, and a controlled test of isolated versus conditional decisions. Its resource axis is protected-bit count, and its measured results come from new frozen executable studies. It neither reuses the earlier hardware performance claims nor presents the general design-flow concept as new.
